%!TEX root = ../main.tex
\chapter{Conclusioni}
\label{cap:conclusioni}

Il Capitolo~\ref{cap:introduzione} ha dichiarato tre obiettivi per questo
lavoro: efficienza, equità e generalizzazione. Questo capitolo li riprende
uno a uno alla luce dei risultati del Capitolo~\ref{cap:confronto}, indica
le direzioni lasciate aperte e discute infine i limiti dello studio.

Il confronto di base (§\ref{sec:risultati-base}) conferma che sia
MaxPressure sia MetaSTGAT Pro battono nettamente Fixed-Time su ogni
metrica, con l'unica eccezione del max wait in training, dove la
rotazione fissa delle fasi di Fixed-Time impone per costruzione un
limite all'attesa che le configurazioni di training, meno dense, non
arrivano a violare: un limite che sparisce non appena il traffico si
fa più intenso in test. Tra MaxPressure e Pro il compromesso al centro di
questo lavoro emerge già a questo livello: MaxPressure resta la scelta più
efficiente in media, Pro riduce il max wait del 34-41\% pur restando entro
pochi punti percentuali su travel time medio.

Lo studio sul peso $\alpha$ quantifica questo compromesso direttamente:
max wait e $tt_{\max}$ peggiorano in modo pressoché monotono al diminuire
di $\alpha$, mentre il travel time medio non peggiora in modo altrettanto
sistematico: $\alpha=0.04$ lo riduce persino sotto quello di MaxPressure
in test, il primo caso in questo lavoro in cui una variante RL supera
MaxPressure su questa metrica in generalizzazione. Il caso più netto a
sostegno della necessità del termine anti-starvation resta però
$\alpha=0$: rimuovere la sola penalità, lasciando intatti stato, maschera
e vincolo di visibilità, è sufficiente da solo a produrre un blocco quasi
totale di un veicolo per l'intera durata dell'episodio su quattro dei
cinque seed della griglia $5\times5$ più densa mai vista in training. È una
prova diretta, non solo correlazionale, che il termine di ricompensa (non
lo stato, non la maschera) è il meccanismo causale dietro la garanzia di
equità osservata. Per questo motivo, nonostante $\alpha=0.04$ sia
leggermente più efficiente in media, questo lavoro adotta $\alpha=0.08$
come riferimento per gli studi successivi: coerentemente con la priorità
dichiarata tra le metriche, è la configurazione più incisiva proprio sulle
due metriche di caso peggiore che questo lavoro tiene in primo piano.

Lo studio di ablation isola un effetto complementare: dei cinque preset
confrontati, solo i due che rimuovono interamente il segnale di attesa
dallo stato o dalla ricompensa (\texttt{paper} ed \texttt{environment})
degradano marcatamente verso il comportamento di MaxPressure, mentre le
tre varianti che disattivano singolarmente meccanismi dell'algoritmo di
apprendimento (\texttt{temporal}, \texttt{single\_dqn},
\texttt{vanilla\_buffer}) restano vicine al riferimento o lo superano su
alcune metriche. Su 35 istanze di test, \texttt{paper} ed
\texttt{environment} sono gli unici, insieme a MaxPressure, a bloccare un
veicolo per quasi l'intera durata dell'episodio; nessuna delle varianti
che conserva stato e ricompensa lo fa mai. La garanzia di equità di questo
lavoro dipende quindi dalla combinazione di segnale di attesa nello stato
e termine di ricompensa dedicato, non da un singolo dettaglio
implementativo dell'algoritmo di apprendimento sottostante.

Il confronto tra meccanismi spaziali aggiunge una terza componente
necessaria: a parità di stato, ricompensa e maschera, MetaSTGCN (l'unico
meccanismo a peso di aggregazione fissato dalla sola topologia, non
appreso per coppia di nodi) resta sistematicamente indietro su ogni
metrica, in un caso peggiore persino di MaxPressure nel proprio stesso
regime di addestramento. Equità ed efficienza non sono quindi solo una
funzione della ricompensa: richiedono anche un meccanismo di
comunicazione capace di instradare il segnale del vicino in difficoltà
fino a chi può agire su di esso.

Ogni modello di questo lavoro è addestrato una sola volta su una griglia
$4\times4$ e valutato, senza riaddestramento, su topologie più grandi
($5\times5$, $6\times6$), su una spaziatura stradale diversa e su profili
di traffico differenti. Il risultato principale è positivo ma non
uniforme. Sulle configurazioni mai viste in training, MetaSTGAT Pro
conserva il proprio vantaggio su MaxPressure e Fixed-Time su tutte le
metriche principali, e le varianti che disattivano meccanismi
dell'algoritmo di apprendimento (\texttt{single\_dqn},
\texttt{vanilla\_buffer}) mostrano un vantaggio sul riferimento
concentrato quasi esclusivamente su queste topologie più grandi: alla
scala di training, lo stesso vantaggio si riduce o si inverte a seconda
della metrica, segno che quei meccanismi aiutano soprattutto la
generalizzazione di scala, non le prestazioni in senso assoluto.

Un secondo asse di generalizzazione, la spaziatura stradale, rivela un
meccanismo distinto: il preset \texttt{environment}, che tra le altre cose
disattiva il vincolo sul campo visivo, degrada in modo sproporzionato
proprio sulle due configurazioni con spaziatura di 200 metri, mai viste in
training a nessuna spaziatura diversa da 100 metri. La causa più diretta è
che quel vincolo limita per costruzione lo scarto tra la distribuzione di
training e quella di test; senza di esso, il modello osserva conteggi di
veicoli mai incontrati durante l'addestramento. Lo stesso pattern,
attenuato, si osserva nel margine di vantaggio di \texttt{single\_dqn} e
\texttt{vanilla\_buffer} sul max wait, che si restringe quasi a scomparire
sulle configurazioni a 200 metri pur restando ampio sulle griglie più
grandi: generalizzare a una scala di rete diversa e generalizzare a una
spaziatura fisica diversa non sono la stessa cosa, e i meccanismi che
aiutano l'una non aiutano necessariamente l'altra.

Il caso di MetaSTSONAR chiude il quadro con un risultato meno atteso: a
parità di ogni altro fattore, $L=2$ ricorrenze di propagazione superano
$L=4$ su $tt_{\max}$ sia in training sia in test, quando il vantaggio
teorico di un campo recettivo più ampio suggerirebbe l'opposto. Poiché lo
svantaggio di $L=4$ è già presente alla scala di training, dove non c'è
alcun cambio di topologia da spiegare, la causa non può essere una
generalizzazione peggiore: è più plausibile un limite legato
all'ottimizzazione (il gradiente deve retropropagare attraverso il doppio
dei passi dell'equazione di propagazione, a parità di procedura di
training) o un effetto di \emph{over-smoothing} già alla scala di
training. Questo lavoro non arriva a isolare quale delle due spiegazioni,
se non entrambe, sia corretta.

I risultati discussi finora aprono la strada a diverse direzioni di
ricerca. Una prima direzione, di natura metodologica, è la ripetizione
indipendente dell'addestramento per ciascuna configurazione: distinguere
un effetto architetturale reale dalla varianza tra run indipendenti,
specialmente per i risultati più sottili di questo lavoro (il vantaggio
di \texttt{single\_dqn} e \texttt{vanilla\_buffer} sulla generalizzazione
di scala, il comportamento di MetaSTSONAR a $L=4$), richiederebbe
ripetizioni non disponibili in questo studio. Una seconda direzione,
sul confronto tra meccanismi spaziali, è un confronto a parametri
pareggiati tra MetaSTGAT, MetaSTGCN e MetaSTSONAR, per isolare la
capacità espressiva del meccanismo di aggregazione dal semplice numero di
parametri disponibili. Una terza direzione, applicativa, è l'estensione
della valutazione di generalizzazione a topologie stradali reali oltre
alle griglie sintetiche di questo lavoro. Infine, i risultati sulla
spaziatura stradale lasciano aperta una domanda diretta: se il divario di
generalizzazione osservato dipenda specificamente dal vincolo sul campo
visivo, come i risultati di \texttt{environment} suggeriscono, o più in
generale da come lo stato rappresenta grandezze fisiche assolute invece
che relative alla scala della rete; una domanda che uno studio dedicato,
con configurazioni intermedie tra le spaziature qui testate, potrebbe
risolvere in modo conclusivo dove questo lavoro si limita a
un'osservazione.

Nonostante i risultati ottenuti, alcuni limiti vanno considerati
nell'applicazione pratica di questo lavoro. Ogni modello, RL o ablation,
deriva da un singolo run di addestramento: la batteria a più seed
introdotta nel protocollo sperimentale ripete la sola valutazione, non
l'addestramento, per cui buona parte dei risultati più sottili discussi
in questo capitolo resta compatibile, oltre che con le spiegazioni
proposte, con la varianza tra run di addestramento indipendenti diversi
dall'unico osservato. Il confronto tra meccanismi spaziali, come già
notato, non pareggia il numero di parametri tra le tre famiglie
confrontate: una parte del divario osservato, specialmente a sfavore di
MetaSTGCN, potrebbe in principio riflettere anche questa differenza, non
solo la capacità espressiva del meccanismo di aggregazione in sé. La
valutazione di generalizzazione resta inoltre confinata a reti sintetiche
a griglia regolare, generate parametricamente, mai a una rete stradale
reale: un limite che questo lavoro condivide con la maggior parte della
letteratura citata nell'introduzione, e che la sola estensione a griglie
più grandi non risolve. Allo stesso modo, \textsc{CityFlow} privilegia la
velocità di simulazione sul realismo microscopico del singolo veicolo
(§\ref{sec:background-tsc}): risultati assoluti su un simulatore più
dettagliato, o su dati reali, potrebbero differire pur lasciando invariate
le conclusioni relative tra i controllori confrontati. Infine, il
confronto tra modelli in questo lavoro non include un test di
significatività statistica in senso classico: la consistenza osservata su
più configurazioni di test e più seed di valutazione rende i risultati
qualitativamente solidi, non un sostituto formale di un test di ipotesi.
